\section{Improvements and Simplification Analysis}

The codebase is lean and well-structured for its scope:
approximately 800 lines across 8 source files, with clear single-responsibility boundaries.
No files are unnecessary, and the module dependency graph is acyclic and logical.

The items below are \emph{opportunities}, not defects --- the code is correct and all
round-trips are lossless.
Each item is labelled with a category and a priority.

\subsection{Critical / High-Priority Improvements}

\subsubsection{Double Encoding in Adaptive Mode}

\textbf{Category: Performance. Priority: High. Effort: Low.}

In \texttt{compress()} (adaptive path), Pass~1 calls \texttt{lzss\_encode\_to\_buffer()}
for every block to obtain the compressed size and determine whether to set
\texttt{meta.compressed = 1}.
Pass~2 then re-encodes the \emph{same} block for actual output.

For a $512 \times 512$ image divided into $16 \times 16$ blocks, there are 1024 blocks,
so the encoder runs 2048 times instead of 1024.
Adaptive mode is therefore roughly twice as slow as necessary.

\textbf{Root cause:} the file format requires all block metadata to precede block data
(so the decompressor knows the layout before reading any data).
A naive single-pass encoder cannot write metadata first without knowing sizes upfront.

\textbf{Fix:} in Pass~1, store the encoded bytes --- not just the size --- alongside
\texttt{meta.compressed}.
One option is to extend \texttt{Block} with a \texttt{vector\textless uint8\_t\textgreater{} encoded\_data}
field populated during Pass~1 and consumed (then freed) during Pass~2.
Another option is a parallel \texttt{vector\textless vector\textless uint8\_t\textgreater\textgreater{}} alongside \texttt{blocks}.

\subsubsection{No Unit Tests}

\textbf{Category: Quality / Maintainability. Priority: High. Effort: Medium.}

There is no unit test infrastructure.
Round-trip correctness is only verified by \texttt{make test}, which exercises the
full pipeline on whole image files.
A bug in \texttt{BitWriter::flush()}, \texttt{lzss\_decode()}, or
\texttt{inverse\_model()} would surface only as a corrupted output file,
not as a targeted test failure pointing directly to the broken function.

\textbf{Fix:} add a \texttt{tests/test\_main.cpp} (or similar) that directly invokes:
\begin{itemize}
  \item \texttt{BitWriter}/\texttt{BitReader} round-trip for edge cases
        (empty, 1 bit, non-multiple of 8 bits).
  \item \texttt{forward\_model} / \texttt{inverse\_model} round-trip on ramp, flat, and random data.
  \item \texttt{scan\_horizontal}/\texttt{unscan\_horizontal} and
        \texttt{scan\_vertical}/\texttt{unscan\_vertical} inverses.
  \item \texttt{lzss\_encode} / \texttt{lzss\_decode} round-trip on all-same, alternating,
        and random byte sequences.
\end{itemize}
No test framework is required; a standalone \texttt{main} with assertions is sufficient.

\subsection{Quality / Correctness Improvements}

\subsubsection{Czech Error Messages}

\textbf{Category: Quality. Priority: Medium. Effort: Low.}

Three user-visible strings use Czech:
\begin{itemize}
  \item \texttt{main.cpp}: \textit{``Nelze otevřít vstupní soubor:''} and
        \textit{``Nelze otevřít výstupní soubor:''}
  \item \texttt{codec.cpp}: \textit{``Neplatný formát souboru (chybné magic bytes)''} and
        \textit{``Velikost vstupu \ldots není dělitelná šířkou''}
\end{itemize}

The assignment requires English \emph{comments}; error messages are not comments.
However, mixing Czech runtime output with an otherwise English-commented codebase
is inconsistent and could confuse future readers.

\textbf{Fix:} translate all four strings to English equivalents such as
\textit{``Cannot open input file:''}, \textit{``Invalid file format (bad magic bytes)''}, etc.

\subsubsection{Width Validation Belongs in Argument Parsing}

\textbf{Category: Correctness / Defensive Programming. Priority: Medium. Effort: Low.}

\texttt{args.width} defaults to $-1$ as a sentinel for ``not provided''.
When passed as an \texttt{int} to \texttt{compress()}, which immediately converts it
to \texttt{uint32\_t}, a value of $-1$ becomes $2^{32}-1 = 4{,}294{,}967{,}295$.
The subsequent check \texttt{pixels.size() \% image\_width != 0} may technically
avoid a crash but produces a misleading error message about image dimensions.

\textbf{Fix:} validate \texttt{width > 0} inside \texttt{parse\_arguments()} (in compression mode)
and exit immediately with a clear message such as \textit{``Error: -w must be a positive integer.''}.
This removes an implicit semantic dependency between \texttt{args.hpp} and \texttt{codec.cpp}
around the $-1$ sentinel.

\subsection{Minor / Low-Priority Improvements}

\subsubsection{scan\_horizontal Is an Unnecessary Heap Allocation}

\textbf{Category: Performance / Simplification. Priority: Medium. Effort: Low.}

\texttt{scan\_horizontal()} returns a copy of its input unchanged (row-major storage is
the native layout).
For a $512 \times 512$ image in static mode this allocates 262\,144 bytes unnecessarily.
In adaptive mode, every block choosing horizontal scan incurs the same cost.

\textbf{Options:}
\begin{enumerate}
  \item Add a branch in \texttt{codec.cpp} to skip the scan call for horizontal direction
        and operate directly on the source pixels.
  \item Make \texttt{scan\_horizontal} a no-op that returns the input by move (if the caller
        passes by value) or provide a separate \texttt{scan\_horizontal\_inplace} that returns
        a \texttt{const\&}.
\end{enumerate}
Option~1 is simpler.
The API symmetry between \texttt{scan\_horizontal} and \texttt{scan\_vertical} is worth
keeping for clarity, but the unnecessary copy in the H-path is easy to eliminate.

\subsubsection{lzss\_encode\_to\_buffer and ostringstream Overhead}

\textbf{Category: Performance. Priority: Low. Effort: Medium.}

\texttt{lzss\_encode\_to\_buffer} wraps a \texttt{std::ostringstream} to collect output
bytes, then copies the internal string buffer into a \texttt{vector\textless uint8\_t\textgreater}.
This involves two separate heap allocations.

Since each block is at most 256 bytes ($16 \times 16$) and the LZSS-encoded output is
similarly bounded, the absolute cost is small.
A \texttt{vector}-backed \texttt{streambuf} would eliminate one allocation, but this
optimisation is unlikely to be noticeable in practice.

\subsubsection{BLOCK\_SIZE Defined in codec.cpp Only}

\textbf{Category: Simplification / Cohesion. Priority: Low. Effort: Low.}

\texttt{BLOCK\_SIZE = 16} and \texttt{BLOCK\_SIZE\_LOG = 4} are defined as
\texttt{static constexpr} in \texttt{codec.cpp}, making them invisible to other modules.
Logically, block size belongs with the block-related types (\texttt{Block},
\texttt{BlockMeta}) in \texttt{scanner.hpp}.

\textbf{Fix:} move both constants to \texttt{scanner.hpp} so any future component
that needs to know the block size can find it in the obvious place.

\subsection{What Should Not Be Changed}

The following aspects of the current design are correct and should be preserved:

\begin{itemize}
  \item \textbf{7-file module structure.}
        \texttt{args}, \texttt{bitio}, \texttt{lzss}, \texttt{model}, \texttt{scanner},
        \texttt{codec}, \texttt{main} each have a single clear responsibility.
        Do not collapse \texttt{scanner} and \texttt{model} into \texttt{codec} ---
        the separation aids readability and testability.

  \item \textbf{LZSS sliding-window algorithm.}
        LZSS is the right algorithm for this problem: it exploits spatial correlation
        in images, requires no stored dictionary, and has a trivially simple decoder.
        There is no reason to switch to LZ78/LZW or a statistical method.

  \item \textbf{Flag-group encoding (8 tokens per flag byte).}
        The current byte-aligned flag packing is a standard LZSS practice.
        It keeps output byte-aligned and is easy to decode with a simple bitmask loop.

  \item \textbf{argparse.hpp inclusion.}
        The third-party argument parsing library is not a compression library and
        does not violate the assignment constraint.
        Replacing it with manual \texttt{getopt} parsing would save a few hundred
        lines of header but add maintenance burden with no benefit.

  \item \textbf{Header-only BitIO.}
        \texttt{BitWriter} and \texttt{BitReader} are short enough to be inline.
        A separate \texttt{bitio.cpp} translation unit would add build complexity
        for no benefit.
\end{itemize}

\subsection{Summary Table}

\begin{center}
\begin{tabular}{r l l l l}
\toprule
\textbf{\#} & \textbf{Item} & \textbf{Category} & \textbf{Priority} & \textbf{Effort} \\
\midrule
1 & Double encoding in adaptive mode     & Performance            & High   & Low \\
2 & No unit tests                        & Quality                & High   & Medium \\
3 & Czech error messages                 & Quality                & Medium & Low \\
4 & Width validation in argument parser  & Correctness            & Medium & Low \\
5 & scan\_horizontal unnecessary copy   & Performance            & Medium & Low \\
6 & ostringstream in encode\_to\_buffer  & Performance            & Low    & Medium \\
7 & BLOCK\_SIZE defined in codec.cpp     & Simplification         & Low    & Low \\
\bottomrule
\end{tabular}
\end{center}
